\section{Transverse curved directions and a sharp fourth-jet correction}
\label{sec:curves}

The cubic coordinate formula survives a nonlinear change of spectral
path. The path contributes an explicitly classical correction, while
the coefficients of $\Omega_3$ and $\kappa$ depend only on its tangent.

Using \eqref{ray:h0}--\eqref{ray:hA}, define
\begin{align}
 K(A,B,C)=\frac1{\Gamma(1+C)}\biggl\{
 &\zeta(A)\zeta(B)
 -\frac{C\Gamma(1-A)\zeta(B-1)}{A+C}
 -\frac{C\Gamma(1-B)\zeta(A-1)}{B+C}\notag\\
 &+C\bigl(h_0+h_A(A+B)+h_C\,C\bigr)\biggr\}.
 \label{curve:K}
\end{align}
Although $K$ is meromorphic in three variables, its restriction to
each of the paths in the following theorem is holomorphic.

\begin{theorem}[Cubic transport along a spectral curve]\label{thm:curve}
Let $A(t),B(t),C(t)$ be analytic near zero, vanish at zero, and have
first derivatives $a,b,c$ satisfying \eqref{ray:admissible}.
Write $\mathscr W(t)=\T(A(t),B(t),C(t))$. Then
\begin{equation}\label{curve:main}
 \mathscr W'''(0)=W_{a,b,c}'''(0)+
 6[t^3]\bigl(K(A(t),B(t),C(t))-K(at,bt,ct)\bigr).
\end{equation}
The correction uses only ordinary Gamma and zeta derivatives and
the Taylor coefficients of $A,B,C$ through degree four. In particular,
the coefficients of $\Omega_3$ and $\kappa$ remain $abc$ and
$-c\Delta/2$.
\end{theorem}

\begin{proof}
Write the quadratic homogeneous part of \eqref{ray:H-Taylor} as
$H_2$. Formula \eqref{ray:decomposition} gives
\[
 \T-K=\frac{C}{\Gamma(1+C)}H_2(A,B,C)
            +O\bigl((|A|+|B|+|C|)^4\bigr).
\]
After substitution of the curve, its coefficient of $t^3$ is
$cH_2(a,b,c)$, exactly the coefficient on the tangent ray. Subtracting
the two restrictions proves \eqref{curve:main}.

The only possible need for one additional path coefficient comes
from $C/(A+C)$ and $C/(B+C)$. Remove their common factor $t$.
The new denominators have nonzero constant terms, so the coefficient
of $t^3$ in each ratio uses the original curves only through degree
four. All the remaining compositions use at most degree three.
This proves the finite-jet assertion and holomorphy of the restrictions.
\end{proof}

\subsection{A correction involving only the lower directional layer}

The correction can be made smaller than the Gamma expression
\eqref{curve:K} suggests. Put $v=(a,b,c)$ and write the path as
\[
 (A(t),B(t),C(t))=tv+t^2\alpha+t^3\beta+t^4\delta+O(t^5).
\]
Let $f_j(v)=W_v^{(j)}(0)$ for $j=0,1,2$, explicitly given by
\eqref{ray:lower}. Denote ordinary differentials in the slope
variables by $D$. Then
\begin{align}
 \mathscr W'''(0)-W_v'''(0)={}&
 6Df_0(v)[\delta]+6D^2f_0(v)[\alpha,\beta]
                  +D^3f_0(v)[\alpha,\alpha,\alpha]\notag\\
 &+6Df_1(v)[\beta]+3D^2f_1(v)[\alpha,\alpha]
                  +3Df_2(v)[\alpha].
 \label{curve:differentials}
\end{align}
In particular the correction belongs to the span of
$1,L,L^2,\zeta'(-1),\zeta''(0),\zeta''(-1)$ with rational
functions of the path coefficients. No third zeta derivative or
additional Gamma constant is needed.

To prove \eqref{curve:differentials}, cancel the common factor $t$
in the two polar ratios in \eqref{ray:decomposition}. The function
$(t,v)\mapsto\T(ta,tb,tc)$ is then jointly holomorphic near
$(0,v)$ for admissible $v$. Its Taylor series is
$\sum_{j\ge0}f_j(v)t^j/j!$. Substitute
$v+\alpha t+\beta t^2+\delta t^3+O(t^4)$ for $v$ and take
the coefficient of $t^3$. The ordinary multivariate Taylor formula
gives exactly \eqref{curve:differentials}.

\begin{proposition}[The fourth jet is genuinely needed]\label{prop:sharpcurve}
Fix admissible $a,b,c$ with $c\ne0$. Replacing the straight path
$(at,bt,ct)$ by $(at+\varepsilon t^4,bt,ct)$ changes its cubic
derivative by exactly
\begin{equation}\label{curve:sharp}
 -\frac{c\varepsilon}{2(a+c)^2}.
\end{equation}
Here $\varepsilon$ is the coefficient of $t^4$, not the fourth derivative.
\end{proposition}
\begin{proof}
Only the ratio $C/(A+C)$ in \eqref{curve:K} changes at degree three.
Its multiplier at the origin is $-\zeta(-1)=1/12$. The ratio changes by
$-c\varepsilon t^3/(a+c)^2+O(t^4)$. Multiplication by $3!=6$
gives \eqref{curve:sharp}.
\end{proof}

For example, the curve $(t+\varepsilon t^4,t,t)$ has cubic
derivative $\Omega_3-\varepsilon/8$, even though its first three
path derivatives agree with those of the diagonal.

This theorem concerns transverse curves. If a path lies in $A+C=0$
or $B+C=0$, or is tangent to one of those divisors, a separate
restriction or normal finite-part prescription is needed. The theorem
does not assign a value by substituting a zero denominator into
\eqref{ray:main}.
